\section{Part 11a-2: Fourier and Taylor cutoffs, the maximal homological equation, the canonical step}\label{bpP:p11a2}
\begingroup\sloppy

\newcommand{\Dd}[3]{\mathcal D(#1,#2;#3)}
\newcommand{\nrm}[2]{\lVert #1\rVert_{#2}}
\newcommand{\cN}{c_N}
\newcommand{\EN}{\varepsilon_N}
\newcommand{\Lam}[1]{\Lambda_{#1}}
\newcommand{\Tr}{\mathsf T}
\newcommand{\cut}{\mathcal T}

This is the second of the five files of Part~11a (split plan \texttt{PLAN-11a.md}; the file is independent of File~11a-1 except for the divisor interface of Lemma~\ref{bpP:p11a:limit-bridge}). It proves, in the frozen conventions (angles of period $1$, characters $e^{2\pi ik\cdot\theta}$, max norms), the \emph{analytic} half of the maximal-case KAM step of R\"ussmann, replacing every external or defective ingredient:
\begin{itemize}
\item \texttt{fourier-cutoff}, \texttt{loss-choice}, \texttt{loss-sum} (Lemmas~\ref{bpP:p11a:fourier-cutoff}, \ref{bpP:p11a:loss-choice}, \ref{bpP:p11a:loss-sum}): an elementary Fourier cutoff in $\lvert k\rvert_1$ with an explicit constant $c_N$, the loss $d_T$ with error $\le\alpha$, and the sum of the losses over a geometric schedule. This replaces Rus01 Thm~7.3 (via Akhiezer) and the repaired Rus05 estimate (deviation D-4).
\item \texttt{taylor-cutoff}, \texttt{cutoff} (Lemma~\ref{bpP:p11a:taylor-cutoff}, Proposition~\ref{bpP:p11a:cutoff}): the Cauchy--Taylor cutoff of order $1$ in the action variable and the composed approximation $R\mapsto R_*$ (replacing Rus01 Thms~7.2, 7.4; deviation D-5).
\item \texttt{param-cauchy}, \texttt{divisor-wide}, \texttt{divisor-perturb}, \texttt{homological-maximal} (Lemmas~\ref{bpP:p11a:param-cauchy}, \ref{bpP:p11a:divisor-wide}, \ref{bpP:p11a:divisor-perturb}, Proposition~\ref{bpP:p11a:homological-maximal}): parameter Cauchy estimates of every order, the cutoff-wide divisor from the per-mode interface of File~11a-1, its extension to a complex neighbourhood of the parameter set, and the maximal homological equation with $\eta$-linear normal form (Rus01 \S11, pp.~148--153, case $p=q=0$).
\item \texttt{affine-flow}, \texttt{real-lipschitz}, \texttt{canonical-step} (Lemma~\ref{bpP:p11a:affine-flow}, Lemma~\ref{bpP:p11a:real-lipschitz}, Proposition~\ref{bpP:p11a:canonical-step}): the time-one flow of $\varphi_0(\theta)+\varphi_1(\theta)\cdot\eta$ as a simple canonical map (Rus01 \S12, pp.~154--155), through track~N's flow package, and the one-step conjugation $H\circ\Phi=N_++R_+$ with every domain loss and constant explicit.
\end{itemize}

\subsection{Conventions, domains, external facts}\label{bpP:p11a:c2-conv}

\paragraph{Variables.} $N\ge1$ is the number of angles and of actions. The point of the extended phase space is $(\theta,\eta,b)\in\C^N\times\C^N\times\C^N$; $\theta$ is an angle (period one), $\eta$ an action \emph{translated to $0$}, $b$ the parameter of the family (Rus01's $b$, \S14). All norms on $\R^N$, $\C^N$ are max norms $\lvert\cdot\rvert_\infty$, and $\lvert k\rvert_1=\sum\lvert k_i\rvert$; for a matrix $\lvert M\rvert_\infty$ is the induced norm (maximal row sum); for a symmetric $j$-linear map $\lVert D^ju\rVert=\sup\{\lvert D^ju(a_1,\dots,a_j)\rvert_\infty:\lvert a_i\rvert_\infty\le1\}$. For $E\subset\R^N$ and $t>0$ put $E+t=\{b\in\C^N:\exists b'\in E,\ \lvert b-b'\rvert_\infty<t\}$, an open set, and for $b'\in E$ the polydisc $P(b',t)\subset E+t$.
\begin{definition}[Domains and sup norms]\label{bpP:p11a:domain-def}
For $r,s>0$ and $B\subset\C^N$ put $T_r=\{\theta\in\C^N:\lvert\operatorname{Im}\theta\rvert_\infty<r\}$ (\Kref{bp:p0:domains}) and
\[\Dd rsB=\{(\theta,\eta,b):\ \theta\in T_r,\ \lvert\eta\rvert_\infty<s,\ b\in B\}.\]
For $F$ on $\Dd rsB$, or on $T_r\times B$ ($s=0$ is not written; the norm is then $\nrm F{r;B}$), $\nrm F{r,s;B}=\sup\lvert F\rvert_\infty$ over the domain. A function is \emph{admissible} if it is holomorphic in all its variables ($B$ open) and $1$-periodic in each $\theta_j$. For $T\ge0$ the \emph{truncation} of a bounded admissible (or merely continuous periodic) $F$ is $\cut_TF=\sum_{\lvert k\rvert_1\le T}\hat F_k\chi_k$, with $\hat F_k=\int_{[0,1]^N}F(x)\chi_{-k}(x)\,dx$ (coefficients taken for each fixed $(\eta,b)$) and $\chi_k=e^{2\pi ik\cdot\theta}$; its average is $\langle F\rangle=\hat F_0$. A function $F$ is a \emph{trigonometric polynomial of degree $\le T$} if $F=\cut_TF$.
\end{definition}
\textbf{Convention for the bracket and the field} (that of track~N, \Nref{bpN:p5:setting}): $X_\phi=(\partial_\eta\phi,-\partial_\theta\phi)$ on $\C^N\times\C^N$ (angles first), $\{u,v\}=Du\cdot X_v=\partial_\theta u\cdot\partial_\eta v-\partial_\eta u\cdot\partial_\theta v$, and the generating function acts by the time-one flow $\Phi=X_\phi^1$. This is the same scalar Hamiltonian field and bracket as in Rus01 after the conversion of rows to columns (Rus01 uses row vectors and $J_{\mathrm{src}}=\bigl(\begin{smallmatrix}0&-\mathbf 1\\ \mathbf 1&0\end{smallmatrix}\bigr)$, so that its row field $S_xJ_{\mathrm{src}}=(S_y,-S_x)$ is the transpose of $X_S$ above); the homological equation below has the same positive sign as Rus01's $\{S,N\}=\partial_\omega S$, (11.12), up to the factor $2\pi$ of the period-one convention. The sign was fixed once to match \Nref{bpN:p5:taylor}; it does not affect any estimate.

\paragraph{External facts.} Besides \Kref{bp:p1:ext-cauchy}, \Kref{bp:p1:ext-mvt}, \Kref{bp:p1:ext-integral} and \Nref{bpN:p4:ext-param}:
\begin{enumerate}[label=(X\arabic*),ref=(X\arabic*)]
\item\label{bpP:p11a:x-cauchy} \emph{Cauchy coefficient estimate in one variable.} If $g$ is holomorphic on $\lvert\zeta\rvert<R$ with values in $\C^q$ (max norm) and $\lvert g\rvert\le M$, then $g(\zeta)=\sum_{j\ge0}g_j\zeta^j$ on $\lvert\zeta\rvert<R$ with $\lvert g_j\rvert\le MR^{-j}$. (Cauchy integral formula for the coefficients; Mathlib: \nolinkurl{Complex.hasSum_taylorSeries_on_ball}, \nolinkurl{norm_cauchyPowerSeries_le}, names not machine-checked here.) Applied successively in several variables, if $G$ is holomorphic and bounded by $M$ on a polydisc of radii $R_1,\dots,R_j$ then $\lvert\partial_{\zeta_1}\cdots\partial_{\zeta_j}G(0)\rvert\le M\prod_iR_i^{-1}$.
\item\label{bpP:p11a:x-duality} $\lvert k\cdot x\rvert\le\lvert k\rvert_1\lvert x\rvert_\infty$ for $k\in\Z^N$, $x\in\C^N$, and $\lvert\chi_k(\theta)\rvert\le e^{2\pi\lvert k\rvert_1\lvert\operatorname{Im}\theta\rvert_\infty}$.
\item\label{bpP:p11a:x-ode} A linear inhomogeneous ODE $y'=P(t)y+q(t)$ with continuous coefficients on $[0,1]$ has a unique global solution, affine in $y(0)$; Gronwall: if $y'=P(t)y$, $\lvert P\rvert\le\Lambda$, then $\lvert y(t)\rvert\le e^{\Lambda t}\lvert y(0)\rvert$. (Mathlib: \nolinkurl{ODE_solution_unique}, Gronwall.)
\end{enumerate}

\paragraph{Lean totality.} Nothing here divides by zero: the only division is by $2\pi i\,k\cdot\omega(b)$ for $0<\lvert k\rvert_1\le T$, where the hypotheses give $\lvert k\cdot\omega(b)\rvert\ge D/2>0$ (Lemma~\ref{bpP:p11a:divisor-perturb}); $\cut_T$ is a finite sum; no $\mathrm{fderiv}$ of a non-differentiable function occurs, because every derivative taken is of an admissible function on its open domain.

%%%%%%%%%%%%%%%%%%%%%%%%%%%%%%%%%%%%%%%%%%%%%%%%%%%%%%%%%%%%%%%%%%%%%%%%%%%%%%%%
\subsection{The elementary Fourier cutoff and the choice of the loss}\label{bpP:p11a:sec-fourier}

\begin{lemma}[Fourier cutoff in $\lvert k\rvert_1$]\label{bpP:p11a:fourier-cutoff}
Let $N\ge1$, $r>0$ and put
\[c_N=2\Bigl(1+\frac1\pi\Bigr)\Bigl(\frac{3N}{e\pi}\Bigr)^{N-1},\qquad \varepsilon_N(T,d)=c_N\,(1+d^{-N})\,e^{-\pi Td}\qquad(T\ge0,\ d>0).\]
Then $c_N\ge1$. Let $F:T_r\to\C^q$ (max norm) be holomorphic, $1$-periodic and bounded, $T\ge0$ and $0<d<r$. Then $\cut_TF$ is a trigonometric polynomial of degree $\le T$, entire in $\theta$, and
\begin{enumerate}[label=(\alph*)]
\item $\nrm{F-\cut_TF}{r-d}\le\varepsilon_N(T,d)\,\nrm Fr$;
\item $\nrm{\cut_TF}{r-d}\le(1+\varepsilon_N(T,d))\,\nrm Fr$;
\item if $F$ is real on $\R^N$, then so is $\cut_TF$.
\end{enumerate}
The dimension $N$ is that of the angle; there is no rescaling of $T$ or $d$.
\source{Ours. Replaces Rus01 Thm~7.3 p.~134 (approximation by trigonometric polynomials with multipliers, from [21] via Akhiezer's theorem) and its repair Rus05 \S\S9--11 pp.~26--31 (deviation D-4). Idea: Cauchy decay of coefficients (\Kref{bp:p1:decay}) and a lattice count (\Kref{bp:p4:lattice-count}); no multipliers, no smoothing.}
\end{lemma}
\begin{proof}
\proofstep{$c_N\ge1$}\label{stepP:p11a:fc-const} $c_1=2(1+1/\pi)>2$. For $N\ge2$, $c_N=2(1+1/\pi)(3N/(e\pi))^{N-1}$ and $3N/(e\pi)\ge6/(e\pi)>0.70$; for $N=2$, $c_2>2\cdot1.31\cdot0.70>1.8$; for $N\ge3$, $3N/(e\pi)\ge9/(e\pi)>1$, so $c_N\ge2$.

\proofstep{Decay and Fourier series}\label{stepP:p11a:fc-decay} By \Kref{bp:p1:decay} (with $\rho'\uparrow r$), $\lvert\hat F_k\rvert\le e^{-2\pi\lvert k\rvert_1r}\nrm Fr$ (componentwise, hence in the max norm). Hence for every $\rho<r$, $\sum_k\lvert\hat F_k\rvert e^{2\pi\lvert k\rvert_1\rho}<\infty$, and by \Kref{bp:p1:summation} $F=\sum_k\hat F_k\chi_k$ on $T_\rho$; $\cut_TF$ is a finite sum, entire. On $T_{r-d}$, $F-\cut_TF=\sum_{\lvert k\rvert_1>T}\hat F_k\chi_k$. For (c): if $F$ is real on $\R^N$ then $\hat F_{-k}=\overline{\hat F_k}$ and $\cut_TF$ is a real sum.

\proofstep{Tail}\label{stepP:p11a:fc-tail} By \ref{bpP:p11a:x-duality}, for $\theta\in T_{r-d}$,
$\lvert F(\theta)-\cut_TF(\theta)\rvert\le\sum_{\lvert k\rvert_1>T}\nrm Fr\,e^{-2\pi\lvert k\rvert_1r}e^{2\pi\lvert k\rvert_1(r-d)}=\nrm Fr\sum_{\lvert k\rvert_1>T}e^{-2\pi\lvert k\rvert_1d}$.

\proofstep{Count}\label{stepP:p11a:fc-count} Group by $j=\lvert k\rvert_1$ ($j\ge1$ integer, $j>T$) and use $\#\{\lvert k\rvert_1=j\}\le2\cdot3^{N-1}j^{N-1}$ (\Kref{bp:p4:lattice-count}):
$\sum_{\lvert k\rvert_1>T}e^{-2\pi\lvert k\rvert_1d}\le2\cdot3^{N-1}\sum_{j>T}j^{N-1}e^{-\pi jd}\,e^{-\pi jd}$.

\proofstep{Polynomial against exponential}\label{stepP:p11a:fc-poly} By \Kref{bp:p1:exp-sums}(a) with $\tau=N-1$ and $a=\pi d$: $j^{N-1}e^{-\pi jd}\le\bigl(\frac{N-1}{e\pi d}\bigr)^{N-1}$ (read $0^0=1$). The integers $j>T$ are $j_0,j_0+1,\dots$ with $j_0>T$, so $\sum_{j>T}e^{-\pi jd}\le e^{-\pi j_0d}/(1-e^{-\pi d})\le e^{-\pi Td}(1+\frac1{\pi d})$, since $1/(1-e^{-x})=1+1/(e^x-1)\le1+1/x$.

\proofstep{Assembly}\label{stepP:p11a:fc-assemble} The tail sum is at most $2\cdot3^{N-1}\bigl(\frac{N-1}{e\pi}\bigr)^{N-1}\bigl(d^{-(N-1)}+\frac{d^{-N}}\pi\bigr)e^{-\pi Td}$. Now $d^{-(N-1)}\le1+d^{-N}$ (if $d\ge1$ the left side is $\le1$, if $d<1$ it is $\le d^{-N}$), so the bracket is $\le(1+\frac1\pi)(1+d^{-N})$; and $N-1\le N$ gives the factor $(3N/(e\pi))^{N-1}$. This is (a); (b) is (a) and the triangle inequality ($\nrm F{r-d}\le\nrm Fr$).
\end{proof}

\begin{lemma}[Choice of the loss $d_T$]\label{bpP:p11a:loss-choice}
Let $N\ge1$, $0<\alpha\le1$, $T\ge1$ and
\[d_T=d_T(\alpha)=\frac{\log(2c_N/\alpha)+N\log(\pi T)}{\pi T}\ \ (>0).\]
If $d_T\le1$, then $\varepsilon_N(T,d_T)\le\alpha$. Moreover $T\mapsto d_T(\alpha)$ is decreasing on $[1,\infty)$ for fixed $\alpha$, and $\alpha\mapsto d_T(\alpha)$ is decreasing in $\alpha$.
\source{Ours; corrects the first draft of \texttt{PLAN-11a.md}, where $d=(N+1)\log(1/\alpha)/(2\pi T)$ does not give error $\le\alpha$ because the factor $d^{-N}$ grows like $T^N$ (review item P11A1-P02).}
\end{lemma}
\begin{proof}
\proofstep{The exponent}\label{stepP:p11a:lc-exponent} Put $A=\pi Td_T=\log(2c_N/\alpha)+N\log(\pi T)$. Since $c_N\ge1$, $\alpha\le1$, $T\ge1$: $A\ge\log2+N\log\pi>1$.

\proofstep{The estimate}\label{stepP:p11a:lc-estimate} For $d=d_T\le1$: $1+d^{-N}\le2d^{-N}$ and $d^{-N}=(\pi T/A)^N\le(\pi T)^N$ ($A\ge1$). And $e^{-\pi Td}=e^{-A}=\frac{\alpha}{2c_N}(\pi T)^{-N}$. Hence $\varepsilon_N(T,d_T)\le c_N\cdot2(\pi T)^N\cdot\frac\alpha{2c_N}(\pi T)^{-N}=\alpha$.

\proofstep{Monotonicity}\label{stepP:p11a:lc-mono} With $a=\log(2c_N/\alpha)\ge\log2>0$: $\frac d{dT}\frac{a+N\log(\pi T)}{\pi T}=\frac{\pi(N-a-N\log(\pi T))}{(\pi T)^2}<0$ because $\pi T\ge\pi>e$. The $\alpha$-monotonicity is clear.
\end{proof}

\begin{lemma}[Sum of the losses on a geometric schedule]\label{bpP:p11a:loss-sum}
Let $\Theta>1$, $T_0\ge1$, $T_\nu=T_0\Theta^\nu$, and let $\alpha_\nu\in(0,1]$ satisfy $\log(2c_N/\alpha_\nu)\le a+b\nu$ ($a,b\ge0$). Then $d_\nu:=d_{T_\nu}(\alpha_\nu)\le(a+b\nu+N\log(\pi T_0)+N\nu\log\Theta)/(\pi T_0\Theta^\nu)$ and
\[\sum_{\nu\ge0}d_\nu\le\frac1{\pi T_0}\Bigl[\bigl(a+N\log(\pi T_0)\bigr)\frac\Theta{\Theta-1}+(b+N\log\Theta)\frac\Theta{(\Theta-1)^2}\Bigr].\]
In particular for $\alpha_\nu=\alpha_0\delta^{\kappa\nu}$ ($0<\delta<1$): $a=\log(2c_N/\alpha_0)$, $b=\kappa\log(1/\delta)$. The hypothesis $d_\nu\le1$ needed in Lemma~\ref{bpP:p11a:loss-choice} is \emph{not} implied and must be checked by the consumer (it holds for all $\nu$ as soon as the displayed bound is $\le1$).
\source{Ours.}
\end{lemma}
\begin{proof}
$\log(\pi T_\nu)=\log(\pi T_0)+\nu\log\Theta$ and the definition of $d_T$ give the bound on $d_\nu$; sum with $\sum_{\nu\ge0}\Theta^{-\nu}=\Theta/(\Theta-1)$ and $\sum_{\nu\ge0}\nu\Theta^{-\nu}=\Theta/(\Theta-1)^2$.
\end{proof}

%%%%%%%%%%%%%%%%%%%%%%%%%%%%%%%%%%%%%%%%%%%%%%%%%%%%%%%%%%%%%%%%%%%%%%%%%%%%%%%%
\subsection{The Taylor cutoff in the action and the composed cutoff}\label{bpP:p11a:sec-taylor}

\begin{lemma}[Taylor cutoff of order $1$ in $\eta$]\label{bpP:p11a:taylor-cutoff}
Let $r,s>0$, $B\subset\C^N$ open, $F:\Dd rsB\to\C^q$ admissible, $M=\nrm F{r,s;B}<\infty$. Put $F^{(\le1)}(\theta,\eta,b)=F(\theta,0,b)+\partial_\eta F(\theta,0,b)\,\eta$ ($\partial_\eta F$ the Jacobian in $\eta$). Then $F^{(\le1)}$ is admissible (entire in $\eta$) and, for $\theta\in T_r$, $b\in B$:
\begin{enumerate}[label=(\alph*)]
\item $\lvert\partial_{\eta_j}F(\theta,0,b)\rvert_\infty\le M/s$ for each $j$; $\lvert F^{(\le1)}(\theta,\eta,b)\rvert_\infty\le M(1+\lvert\eta\rvert_\infty/s)$, so $\nrm{F^{(\le1)}}{r,s;B}\le2M$;
\item for $0<q<1$: $\nrm{F-F^{(\le1)}}{r,qs;B}\le M\,\dfrac{q^2}{1-q}$.
\end{enumerate}
\source{Ours; replaces Rus01 Thm~7.4 p.~135 (algebraic polynomials, from Lemma~6 of [18], external) in the case needed here, polynomials of degree $\le1$ (deviation D-5). The $q$ plays the role of $\delta^{\beta}$ of Rus01 (7.2).}
\end{lemma}
\begin{proof}
\proofstep{One variable}\label{stepP:p11a:tc-line} Fix $(\theta,b)$ and $\eta\ne0$ with $a=\lvert\eta\rvert_\infty<s$ ($\eta=0$ is trivial). $g(\zeta)=F(\theta,\zeta\eta/a,b)$ is holomorphic on $\lvert\zeta\rvert<s$ (since $\lvert\zeta\eta/a\rvert_\infty=\lvert\zeta\rvert$) and bounded by $M$. By \ref{bpP:p11a:x-cauchy}, $g(\zeta)=\sum g_j\zeta^j$ with $\lvert g_j\rvert\le Ms^{-j}$, and $g_0=F(\theta,0,b)$, $g_1=\partial_\eta F(\theta,0,b)\eta/a$.

\proofstep{(a)}\label{stepP:p11a:tc-a} $\lvert F^{(\le1)}\rvert=\lvert g_0+g_1a\rvert\le M+Ma/s$. For the column bound apply the same argument along $e_j$ ($a=s^-$): $\lvert\partial_{\eta_j}F(\theta,0,b)\rvert\le M/s$ (\Kref{bp:p1:ext-cauchy}).

\proofstep{(b)}\label{stepP:p11a:tc-b} At $\zeta=a$: $\lvert F(\theta,\eta,b)-F^{(\le1)}(\theta,\eta,b)\rvert=\lvert\sum_{j\ge2}g_ja^j\rvert\le M\sum_{j\ge2}(a/s)^j=M\frac{(a/s)^2}{1-a/s}$, increasing in $a$; for $a<qs$ it is $\le Mq^2/(1-q)$.
\end{proof}

\begin{proposition}[The composed cutoff $R\mapsto R_*$]\label{bpP:p11a:cutoff}
Let $r>d>0$, $s>0$, $B\subset\C^N$ open, $R:\Dd rsB\to\C$ admissible with $\varepsilon=\nrm R{r,s;B}<\infty$; $T\ge0$, $0<q<1$, and $E_T:=\varepsilon_N(T,d)$ (Lemma~\ref{bpP:p11a:fourier-cutoff}). Put
\[R_*=\cut_T\bigl(R^{(\le1)}\bigr)=R_{*0}(\theta,b)+R_{*1}(\theta,b)\cdot\eta,\quad R_{*0}=\cut_T[R(\cdot,0,\cdot)],\ \ R_{*1,j}=\cut_T[\partial_{\eta_j}R(\cdot,0,\cdot)].\]
Then $R_{*0}$, $R_{*1,j}$ are trigonometric polynomials of degree $\le T$ in $\theta$, admissible on $\C^N\times B$ (holomorphic in $(\theta,b)$, $1$-periodic), and
\begin{enumerate}[label=(\alph*)]
\item $\nrm{R_{*0}}{r-d;B}\le(1+E_T)\varepsilon$, $\ \nrm{R_{*1,j}}{r-d;B}\le(1+E_T)\varepsilon/s$ for each $j$, $\ \nrm{R_*}{r-d,s;B}\le2(1+E_T)\varepsilon$;
\item $\nrm{R-R_*}{r-d,qs;B}\le\varepsilon\bigl(\frac{q^2}{1-q}+2E_T\bigr)$;
\item if $R$ is real for real $(\theta,\eta,b)$, so is $R_*$.
\end{enumerate}
With $T=T_\nu$, $d=d_\nu$, $\alpha=\alpha_\nu$ of Lemma~\ref{bpP:p11a:loss-choice}: $E_T\le\alpha$ and (b) reads $\le\varepsilon(q^2/(1-q)+2\alpha)$.
\source{Rus01 Thm~7.2 p.~134 and (13.10) p.~159 (the composed estimate $\lvert F-F_{T\beta\delta}\rvert\le2\delta\lvert F\rvert$); the statement and the constants are ours (D-4, D-5).}
\end{proposition}
\begin{proof}
\proofstep{Coefficients}\label{stepP:p11a:cu-coeff} Holomorphy of $R_{*0},R_{*1,j}$ in $(\theta,b)$: the Fourier coefficients $b\mapsto\int_{[0,1]^N}R(x,0,b)\chi_{-k}(x)\,dx$ are holomorphic by \Nref{bpN:p4:ext-param} (the integrand is jointly continuous with continuous $b$-derivative, $R$ being holomorphic on the open set), and $R_*$ is a finite sum of them times characters.

\proofstep{(a)}\label{stepP:p11a:cu-a} Apply Lemma~\ref{bpP:p11a:fourier-cutoff}(b) for each fixed $(\eta,b)$ to $F=R^{(\le1)}(\cdot,\eta,b)$ on $T_r$ with $\nrm Fr\le2\varepsilon$ (Lemma~\ref{bpP:p11a:taylor-cutoff}(a)): $\nrm{R_*(\cdot,\eta,b)}{r-d}\le2(1+E_T)\varepsilon$. For $R_{*0}$ take $F=R(\cdot,0,b)$ ($\nrm Fr\le\varepsilon$) and for $R_{*1,j}$ take $F=\partial_{\eta_j}R(\cdot,0,b)$ ($\nrm Fr\le\varepsilon/s$, Lemma~\ref{bpP:p11a:taylor-cutoff}(a)).

\proofstep{(b)}\label{stepP:p11a:cu-b} $R-R_*=(R-R^{(\le1)})+(R^{(\le1)}-\cut_TR^{(\le1)})$. On $\Dd{r-d}{qs}B\subset\Dd r{qs}B$ the first term is $\le\varepsilon q^2/(1-q)$ (Lemma~\ref{bpP:p11a:taylor-cutoff}(b)); the second is $\le E_T\cdot2\varepsilon$ by Lemma~\ref{bpP:p11a:fourier-cutoff}(a) applied for each $(\eta,b)$.

\proofstep{(c)}\label{stepP:p11a:cu-c} $R^{(\le1)}$ is real at real points ($\partial_\eta R$ is real there); then Lemma~\ref{bpP:p11a:fourier-cutoff}(c).
\end{proof}

%%%%%%%%%%%%%%%%%%%%%%%%%%%%%%%%%%%%%%%%%%%%%%%%%%%%%%%%%%%%%%%%%%%%%%%%%%%%%%%%
\subsection{Parameter estimates and divisors}\label{bpP:p11a:sec-param}

\begin{lemma}[Cauchy estimates in the parameter, every order]\label{bpP:p11a:param-cauchy}
Let $E\subset\R^N$, $0<t'<t$, and $u:E+t\to\C^q$ holomorphic with $\lvert u\rvert_\infty\le M$. Then for every $b\in E+t'$ and $j\ge1$: $\lVert D^ju(b)\rVert\le j^j M\,(t-t')^{-j}$ (norm as in the conventions). The same holds for a function of $(\theta,\eta,b)$ on $\Dd rs{E+t}$ bounded by $M$, for the $b$-derivatives, at every $(\theta,\eta)$.
\source{Ours (the general-order version; \Kref{bp:p1:cauchy-open} covers only orders $\le3$, \Nref{bpN:p4:general-cauchy} only the first derivative).}
\end{lemma}
\begin{proof}
\proofstep{Polydisc}\label{stepP:p11a:pc-poly} Put $\varrho=t-t'$, fix $b\in E+t'$ and $a_1,\dots,a_j$ with $\lvert a_i\rvert_\infty\le1$. There is $b'\in E$ with $\lvert b-b'\rvert<t'$. For $\zeta\in\C^j$ with $\lvert\zeta_i\rvert<\varrho/j$ we have $\lvert\sum\zeta_ia_i\rvert_\infty<\varrho$, so $b+\sum\zeta_ia_i\in P(b',t)\subset E+t$ and $G(\zeta)=u(b+\sum\zeta_ia_i)$ is holomorphic and bounded by $M$ on this polydisc.

\proofstep{Iterated Cauchy}\label{stepP:p11a:pc-iter} $D^ju(b)(a_1,\dots,a_j)=\partial_{\zeta_1}\cdots\partial_{\zeta_j}G(0)$, so by \ref{bpP:p11a:x-cauchy} it is bounded by $M(j/\varrho)^j$. For the vector-valued case take each component and the max norm.
\end{proof}

\begin{lemma}[Cutoff-wide divisor from the per-mode interface]\label{bpP:p11a:divisor-wide}
Let $\gamma>0$, $\tau\ge0$, $T>0$, $v\in\R^N$, and for each $k\in\Z^N$ with $0<\lvert k\rvert_1\le T$ let $u_k\in\R^N$ satisfy $\lvert k\cdot u_k\rvert\ge2\gamma\lvert k\rvert_1^{-\tau}$ and $\lvert u_k-v\rvert_\infty\le\gamma\lvert k\rvert_1^{-\tau-1}$. Then for all $0<\lvert k\rvert_1\le T$,
\[\lvert k\cdot v\rvert\ \ge\ \gamma\lvert k\rvert_1^{-\tau}\ \ge\ \gamma\,T^{-\tau}.\]
\source{Ours; Lemma~\ref{bpP:p11a:limit-bridge} of File~11a-1 and the monotonicity of $x\mapsto x^{-\tau}$.}
\end{lemma}
\begin{proof}
Lemma~\ref{bpP:p11a:limit-bridge} gives the first inequality for each $k$ with $u=u_k$; the second is $\lvert k\rvert_1\le T$, $\tau\ge0$. In the consumer, $u_k=\omega_{\nu(k)}(b')=\tilde\omega_k(b')$ and $v=\omega_\nu(b')$ at stage $\nu$ (File~11a-3, requirements \ref{bpP:p11a:I1}--\ref{bpP:p11a:I3} of Remark~\ref{bpP:p11a:consumers}, with \ref{bpP:p11a:I3} in the corrected form of Remark~\ref{bpP:p11a:i3-corrected}); the hypotheses here are those requirements, evaluated at one parameter point. The \emph{finite-stage} version that the iteration actually proves, for every $b$ in the retained set $K_{\nu+1}$ and without any limit frequency, is Lemma~\ref{bpP:p11a:finite-divisor} (review P11A2-01): its telescoping $\lvert\omega_\nu(b)-\omega_{\nu(k)}(b)\rvert_\infty\le\gamma\lvert k\rvert_1^{-\tau-1}$ is exactly the hypothesis $\lvert u_k-v\rvert\le\gamma\lvert k\rvert_1^{-\tau-1}$ of this lemma.
\end{proof}

\begin{lemma}[Divisors on a complex neighbourhood of the retained set]\label{bpP:p11a:divisor-perturb}
Let $E\subset\R^N$, $t>0$, $T>0$, $D>0$, $\omega:E+t\to\C^N$ holomorphic with $\lVert D\omega(b)\rVert\le L$ for all $b\in E+t$ (operator norm for $\lvert\cdot\rvert_\infty$), and
\[tTL\le D/2,\qquad \lvert k\cdot\omega(b')\rvert\ge D\quad\text{for all }b'\in E,\ 0<\lvert k\rvert_1\le T.\]
Then $\lvert k\cdot\omega(b)\rvert\ge D/2$ for all $b\in E+t$ and $0<\lvert k\rvert_1\le T$.
\source{Rus01 \S11.4 p.~150, step~3 (extension of 11.3.7 from $K_+$ to $B_+$); ours in the maximal case with the power-law $D$ and the max norm.}
\end{lemma}
\begin{proof}
Let $b\in E+t$, $b'\in E$ with $\lvert b-b'\rvert_\infty<t$. The polydisc $P(b',t)\subset E+t$ is convex, so the mean value inequality (\Kref{bp:p1:ext-mvt}, complex normed spaces) gives $\lvert\omega(b)-\omega(b')\rvert_\infty\le Lt$. By \ref{bpP:p11a:x-duality}, $\lvert k\cdot(\omega(b)-\omega(b'))\rvert\le TLt\le D/2$, whence $\lvert k\cdot\omega(b)\rvert\ge D-D/2$.
\end{proof}

%%%%%%%%%%%%%%%%%%%%%%%%%%%%%%%%%%%%%%%%%%%%%%%%%%%%%%%%%%%%%%%%%%%%%%%%%%%%%%%%
\subsection{The maximal homological equation}\label{bpP:p11a:sec-hom}

For $\delta>0$ put $\Lam\delta=\coth(\pi\delta)^N-1=\sum_{k\ne0}e^{-2\pi\delta\lvert k\rvert_1}$ (\Kref{bp:p1:exp-sums}(b) with $x=2\pi\delta$); $\Lam\delta\le(1+\frac1{\pi\delta})^N-1$ and $\Lam\delta\le(4N/(\pi\delta))^N$ (\Kref{bp:p1:exp-sums}(b),(c)).

\begin{proposition}[Homological equation, maximal case with $\eta$-linear normal form]\label{bpP:p11a:homological-maximal}
Let $E\subset\R^N$, $t>0$, $B=E+t$, $T>0$, $D>0$, $\omega:B\to\C^N$ holomorphic, real on $B\cap\R^N$, with $\lVert D\omega\rVert\le L$ on $B$ and the hypotheses of Lemma~\ref{bpP:p11a:divisor-perturb} ($tTL\le D/2$ and $\lvert k\cdot\omega\rvert\ge D$ on $E$ for $0<\lvert k\rvert_1\le T$). Let $r_*>\delta>0$ and let $P_0:\C^N\times B\to\C$, $P_1=(P_{1,j}):\C^N\times B\to\C^N$ be admissible trigonometric polynomials of degree $\le T$ in $\theta$ with $\nrm{P_0}{r_*;B}\le M_0$, $\nrm{P_{1,j}}{r_*;B}\le M_1$. Then there are unique admissible trigonometric polynomials $\varphi_0,\varphi_{1,j}$ of degree $\le T$ with zero average, holomorphic on $\C^N\times B$, such that for every $b\in B$
\[\omega(b)\cdot\partial_\theta\varphi_0=P_0-\langle P_0\rangle,\qquad \omega(b)\cdot\partial_\theta\varphi_{1,j}=P_{1,j}-\langle P_{1,j}\rangle\quad(1\le j\le N),\]
and, with $\Lam\delta$ as above:
\begin{enumerate}[label=(\alph*)]
\item $\nrm{\varphi_0}{r_*-\delta;B}\le\dfrac{\Lam\delta}{\pi D}M_0$, $\ \nrm{\varphi_{1,j}}{r_*-\delta;B}\le\dfrac{\Lam\delta}{\pi D}M_1$;
\item the averages $\langle P_0\rangle$, $\langle P_{1,j}\rangle$ are holomorphic on $B$ and bounded by $M_0$, $M_1$ (no divisor is needed for them);
\item for $0<t'<t$, $j\ge1$, $b\in E+t'$ and $\theta\in T_{r_*-\delta}$: $\lVert D_b^j\varphi_0\rVert\le j^j\frac{\Lam\delta M_0}{\pi D}(t-t')^{-j}$, likewise for $\varphi_{1,j}$ with $M_1$, and for the averages with $M_0,M_1$;
\item if $P_0,P_1$ are real for real $(\theta,b)$, then $\varphi_0,\varphi_1$ and the averages are real for real $(\theta,b)$, $b\in B\cap\R^N$;
\item $\lvert\hat{\varphi}_{1,j,k}(b)\rvert\le M_1e^{-2\pi\lvert k\rvert_1r_*}/(\pi D)$ for $0<\lvert k\rvert_1\le T$, $b\in B$ (used by Lemma~\ref{bpP:p11a:real-lipschitz}).
\end{enumerate}
In the conventions of Section~\ref{bpP:p11a:c2-conv}, for $\varphi=\varphi_0+\varphi_1\cdot\eta$ and $N_b=h+\omega\cdot\eta$: $\{N_b,\varphi\}=-(P-\langle P\rangle)$ with $P=P_0+P_1\cdot\eta$, $\langle P\rangle=\langle P_0\rangle+\langle P_1\rangle\cdot\eta$.
\source{Rus01 \S11.4 pp.~149--151 (equation (11.12), $U_k=G_k/(i\langle k,\omega\rangle)$, and the estimate (11.18), (11.20)--(11.21) with $p=q=0$); sup-norm Fourier estimate (11.2) p.~148 replaced by \Kref{bp:p1:decay}. The cutoff-wide divisor, the loss $\delta$ through $\Lam\delta$ instead of the count $(2T+1)^n$ of (11.21), and the parameter Cauchy estimates are ours.}
\end{proposition}
\begin{proof}
\proofstep{Divisors}\label{stepP:p11a:hm-div} By Lemma~\ref{bpP:p11a:divisor-perturb}, $\lvert k\cdot\omega(b)\rvert\ge D/2$ for all $b\in B$, $0<\lvert k\rvert_1\le T$.

\proofstep{Coefficients}\label{stepP:p11a:hm-coeff} Write $P_0=\sum_{\lvert k\rvert_1\le T}p_k(b)\chi_k$ with $p_k$ holomorphic on $B$ (Step~\ref{stepP:p11a:cu-coeff}). Since $\partial_\omega\chi_k=2\pi i(k\cdot\omega)\chi_k$ (\Kref{bp:p1:fourier-basic}(a)), the equation holds iff $\hat\varphi_{0,k}=p_k/(2\pi i\,k\cdot\omega(b))$ for $0<\lvert k\rvert_1\le T$ and $\hat\varphi_{0,0}=0$; this defines $\varphi_0$ uniquely, holomorphic on $B$ (finite sum of quotients of holomorphic functions with non-vanishing denominator); likewise $\varphi_{1,j}$. Holomorphy in $\theta$ and periodicity are clear.

\proofstep{(a),(e)}\label{stepP:p11a:hm-a} By \Kref{bp:p1:decay}, $\lvert p_k\rvert\le M_0e^{-2\pi\lvert k\rvert_1r_*}$, so $\lvert\hat\varphi_{0,k}\rvert\le M_0e^{-2\pi\lvert k\rvert_1r_*}/(\pi D)$, and for $\theta\in T_{r_*-\delta}$, $\lvert\varphi_0\rvert\le\sum_{k\ne0}\lvert\hat\varphi_{0,k}\rvert e^{2\pi\lvert k\rvert_1(r_*-\delta)}\le\frac{M_0}{\pi D}\sum_{k\ne0}e^{-2\pi\delta\lvert k\rvert_1}=\frac{\Lam\delta M_0}{\pi D}$. Same for $\varphi_{1,j}$; (e) is the coefficient bound just used.

\proofstep{(b),(c)}\label{stepP:p11a:hm-bc} $\langle P_0\rangle=p_0(b)$ is holomorphic, $\lvert p_0\rvert\le M_0$ by \Kref{bp:p1:decay} ($k=0$). (c) is Lemma~\ref{bpP:p11a:param-cauchy} applied, for fixed $\theta$, to $b\mapsto\varphi_0(\theta,b)$ on $B=E+t$ with $M=\Lam\delta M_0/(\pi D)$.

\proofstep{(d)}\label{stepP:p11a:hm-d} At real $b$, $\omega$ is real and $\overline{p_{k}}=p_{-k}$ ($P_0$ real), so $\overline{\hat\varphi_{0,k}}=\overline{p_k}/\overline{2\pi i\,k\cdot\omega}=p_{-k}/(2\pi i(-k)\cdot\omega)=\hat\varphi_{0,-k}$.

\proofstep{The bracket}\label{stepP:p11a:hm-bracket} $\partial_\theta N=0$, $\partial_\eta N=\omega$, so $\{N,\varphi\}=-\omega\cdot\partial_\theta\varphi=-(\omega\cdot\partial_\theta\varphi_0)-(\omega\cdot\partial_\theta\varphi_1)\cdot\eta=-(P-\langle P\rangle)$.
\end{proof}

%%%%%%%%%%%%%%%%%%%%%%%%%%%%%%%%%%%%%%%%%%%%%%%%%%%%%%%%%%%%%%%%%%%%%%%%%%%%%%%%
\subsection{The time-one flow as a simple canonical map}\label{bpP:p11a:sec-flow}

\begin{lemma}[Flow of $\varphi_0(\theta)+\varphi_1(\theta)\cdot\eta$]\label{bpP:p11a:affine-flow}
Let $r_1>2\sigma>0$, $s>0$, $A_1\ge0$ with $4(N+1)A_1\le\sigma$ ($A_1=0$ is allowed: then $\varphi=0$, $\Phi=\mathrm{id}$ and every statement is trivial; the proof below divides by $A_1$ nowhere), and let $\varphi_0:T_{r_1}\to\C$, $\varphi_1:T_{r_1}\to\C^N$ be holomorphic, $1$-periodic, with $\nrm{\varphi_0}{r_1}\le sA_1$ and $\lvert\varphi_{1,j}\rvert\le A_1$ on $T_{r_1}$. Let $\varphi(\theta,\eta)=\varphi_0(\theta)+\varphi_1(\theta)\cdot\eta$, $X_\varphi=(\varphi_1(\theta),\,-\partial_\theta\varphi_0(\theta)-(D\varphi_1(\theta))^{\top}\eta)$, $\Omega=T_{r_1-\sigma}\times B_s$ ($B_s=\{\lvert\eta\rvert_\infty<s\}$), and
\[\kappa=\frac{2(N+1)A_1}{\sigma}\ \le\ \frac12,\qquad A_\eta=\frac{(N+1)sA_1}{\sigma}\ \le\ \frac s4,\qquad \nu=\nu_{s/2,\sigma},\]
where $\nu(a,c)=\max(\lvert a\rvert/\sigma,\ \lvert c\rvert/(s/2))$ (\Nref{bpN:p5:setting}). Put $\Omega_{\mathrm h}=T_{r_1-3\sigma/2}\times B_{3s/4}$ and $\Omega_{\mathrm f}=T_{r_1-2\sigma}\times B_{s/2}$. Then:
\begin{enumerate}[label=(\alph*)]
\item (field) on $\Omega$: $\lvert\partial_\eta\varphi\rvert_\infty\le A_1$, $\lvert\partial_\theta\varphi\rvert_\infty\le A_\eta$, and $\nu(X_\varphi)\le\kappa\le\frac12$.
\item (existence, displacement) $\Phi=X^1_\varphi$ is defined and holomorphic on $\Omega_{\mathrm h}$ ($\Omega_{\mathrm h}\subset\Omega^{1/2}$ in the notation of \Nref{bpN:p5:flow}), the trajectories $X^t_\varphi(z)$, $\lvert t\rvert\le1$, stay in $\Omega$, and $\lvert\Phi_\theta-\theta\rvert_\infty\le A_1$, $\lvert\Phi_\eta-\eta\rvert_\infty\le A_\eta$.
\item (simple form) $\Phi(\theta,\eta)=(\Theta(\theta),\ \mathbf M(\theta)\eta+m(\theta))$ on $\Omega_{\mathrm h}$, where $\Theta$ is the time-one map of the autonomous system $\dot\theta=\varphi_1(\theta)$, and $\mathbf M(\theta)=(D\Theta(\theta))^{-\top}$.
\item (derivative) on $\Omega_{\mathrm f}\subset\Omega^1$: $\lvert D\Phi-\mathbf 1\rvert_\nu\le2\kappa$; in particular $\lvert D\Theta-\mathbf 1\rvert_\infty\le2\kappa$ and $\lvert\mathbf M-\mathbf 1\rvert_\infty\le2\kappa$, and $\lvert\partial_\theta(\mathbf M\eta+m)\rvert_\infty\le\kappa s/\sigma$.
\item (canonical, periodic, real, inverse) $(D\Phi)^\top J\,D\Phi=J$; $\Phi(\theta+k,\eta)=\Phi(\theta,\eta)+(k,0)$ ($k\in\Z^N$); if $\varphi_0,\varphi_1$ are real on $\R^N$, $\Phi$ maps real points of $\Omega_{\mathrm h}$ to real points; $\Psi=X^{-1}_\varphi$ is defined on $\Omega_{\mathrm h}$ and $\Psi\circ\Phi=\mathrm{id}=\Phi\circ\Psi$ on $\Omega_{\mathrm f}$.
\item (parameter) If $\varphi_0,\varphi_1$ depend on $b\in B$ ($B\subset\C^N$ open) so that $(\theta,b)\mapsto\varphi_i(\theta,b)$ are holomorphic with the same bounds for every $b$, then $(b,\theta,\eta)\mapsto\Phi_b(\theta,\eta)$ is holomorphic on $B\times\Omega_{\mathrm h}$, and the same bounds hold for every $b$.
\end{enumerate}
\source{Rus01 \S12 pp.~154--155 (Lemma 12.1, Definition 12.2: simple canonical transformations) for the form; the existence, displacement, derivative, symplectic and reality statements are \Nref{bpN:p5:flow}, \Nref{bpN:p5:symplectic} with $\varepsilon=\kappa$ (Pos93 pp.~211--212 through track~N); new here: the sup-norm field bounds without weighted norms, the simple form, the parameter.}
\end{lemma}
\begin{proof}
\proofstep{(a) Field}\label{stepP:p11a:af-field} By \Kref{bp:p1:cauchy-strip} (with $\rho=r_1$, loss $\sigma$): on $T_{r_1-\sigma}$, $\lvert\partial_{\theta_i}\varphi_0\rvert\le sA_1/\sigma$ and $\lvert\partial_{\theta_i}\varphi_{1,j}\rvert\le A_1/\sigma$. Hence for $\lvert\eta\rvert_\infty<s$: $\lvert\partial_{\theta_i}\varphi\rvert\le sA_1/\sigma+N\cdot A_1 s/\sigma=A_\eta$ ($N$ terms $\lvert\partial_{\theta_i}\varphi_{1,j}\rvert\lvert\eta_j\rvert$), and $\partial_\eta\varphi=\varphi_1$, $\lvert\varphi_1\rvert\le A_1$. Then $\nu(X_\varphi)\le\max(A_1/\sigma,\ A_\eta/(s/2))=\max(A_1/\sigma,2(N+1)A_1/\sigma)=\kappa\le\frac12$; $A_\eta\le s/4$ is $4(N+1)A_1\le\sigma$.

\proofstep{(b) Flow}\label{stepP:p11a:af-flow} $\Omega$ is open and $X_\varphi$ is holomorphic on $\Omega$ with $\nu(X_\varphi)\le\kappa\le\frac12$. \Nref{bpN:p5:flow}(a),(b) give $\Phi$ on $\Omega^{1/2}=\{z:\bar B_\nu(z,\frac12)\subset\Omega\}$, with trajectories in $\Omega$. The ball $\bar B_\nu(z,c)$ is the polydisc $\bar P(z;c\sigma,cs/2)$ (\Nref{bpN:p5:setting}), so $\Omega_{\mathrm h}\subset\Omega^{1/2}$ and $\Omega_{\mathrm f}\subset\Omega^1$. Component bounds: along a trajectory $\dot\theta=\varphi_1$, $\dot\eta=-\partial_\theta\varphi$, bounded by $A_1$, $A_\eta$ (\Kref{bp:p1:ext-mvt}), so $\lvert\Phi_\theta-\theta\rvert\le A_1$, $\lvert\Phi_\eta-\eta\rvert\le A_\eta$.

\proofstep{(c) Simple form}\label{stepP:p11a:af-simple} The $\theta$-equation $\dot\theta=\varphi_1(\theta)$ does not involve $\eta$, so by uniqueness (\Nref{bpN:p5:flow}(a)) the $\theta$-component of the trajectory through $(\theta_0,\eta_0)$ is the solution of this autonomous system through $\theta_0$: $\Phi_\theta(\theta_0,\eta_0)=\Theta(\theta_0)$. Given $\theta(t)$ the $\eta$-equation is linear inhomogeneous, $\dot\eta=-\partial_\theta\varphi_0(\theta(t))-(D\varphi_1(\theta(t)))^\top\eta$, so by \ref{bpP:p11a:x-ode} $\Phi_\eta=\mathbf M(\theta_0)\eta_0+m(\theta_0)$. By \Nref{bpN:p5:symplectic}, $Y=D\Phi=\begin{pmatrix}D\Theta&0\\ *&\mathbf M\end{pmatrix}$ satisfies $Y^\top JY=J$; the $(1,2)$-block of $Y^\top JY$ is $(D\Theta)^\top\mathbf M$ and that of $J$ is $\mathbf 1$, so $\mathbf M=(D\Theta)^{-\top}$.

\proofstep{(d)}\label{stepP:p11a:af-deriv} \Nref{bpN:p5:flow}(c): $\lvert D\Phi-\mathbf 1\rvert_\nu\le2\kappa$ on $\Omega^1$, i.e.\ $\lvert W(D\Phi-\mathbf 1)W^{-1}\rvert_\infty\le2\kappa$ with $W=\operatorname{diag}(\sigma^{-1}\mathbf 1,(s/2)^{-1}\mathbf 1)$. The diagonal blocks of $W(\cdot)W^{-1}$ equal the diagonal blocks of the matrix, and a block has norm at most the row norm of the whole matrix: $\lvert D\Theta-\mathbf 1\rvert,\lvert\mathbf M-\mathbf 1\rvert\le2\kappa$. The lower-left block of $W(D\Phi-\mathbf 1)W^{-1}$ is $(2/s)\sigma\,\partial_\theta(\mathbf M\eta+m)$, whence the last bound ($2\kappa\cdot s/(2\sigma)$).

\proofstep{(e),(f)}\label{stepP:p11a:af-rest} Symplecticity: \Nref{bpN:p5:symplectic}; periodicity and reality: \Nref{bpN:p5:flow}(d),(e) ($\varphi_i$ are $\Z^N$-periodic, so $X_\varphi$ is invariant under $(\theta,\eta)\mapsto(\theta+k,\eta)$); inverse: \Nref{bpN:p5:flow}(c) and $\Psi=X^1_{-\varphi}$. For (f) consider the field $\tilde F(b,z)=(0,X_{\varphi_b}(z))$ on $B\times\Omega$ with the complex norm $\tilde\nu(\beta,z)=\max(\lvert\beta\rvert_\infty/\varsigma,\nu(z))$ for a $\varsigma>0$; it is holomorphic with $\tilde\nu(\tilde F)=\nu(X_{\varphi_b})\le\frac12$, and its flow is $(b,\Phi_b(z))$. By \Nref{bpN:p5:flow}(b) the flow is holomorphic on $\tilde\Omega^{1/2}\supset\{b:P(b,\varsigma/2)\subset B\}\times\Omega^{1/2}$; as $B$ is open every $b$ is covered for $\varsigma$ small.
\end{proof}

\begin{lemma}[Lipschitz bound for the angle map on the reals]\label{bpP:p11a:real-lipschitz}
Let $\varphi_1:\R^N\to\R^N$ be a real trigonometric polynomial of degree $\le T$, $1$-periodic, and $\Theta^t$ the flow of $\dot\theta=\varphi_1(\theta)$ on $\R^N$ (global). Put $\Lambda=\sup_{x\in\R^N}\lvert D\varphi_1(x)\rvert_\infty$. Then for all $\lvert t\rvert\le1$ and $\theta\in\R^N$:
$\lvert D\Theta^t(\theta)-\mathbf 1\rvert_\infty\le e^{\Lambda}-1$, $\ \Theta^t(\theta+k)=\Theta^t(\theta)+k$, $\ \lvert\Theta^t(\theta)-\theta\rvert_\infty\le\sup\lvert\varphi_1\rvert_\infty$, and
\[\Lambda\ \le\ 2\pi T\max_i\sum_{0<\lvert k\rvert_1\le T}\lvert\hat\varphi_{1,i,k}\rvert .\]
If $\lvert\hat\varphi_{1,i,k}\rvert\le M_1e^{-2\pi\lvert k\rvert_1r_*}/(\pi D)$ (Proposition~\ref{bpP:p11a:homological-maximal}(e)), then $\Lambda\le\Lambda_\mathbb R:=2TM_1\Lam{r_*}/D$.
\source{Ours; this is the source of the factor $e^{cC_\nu T_\nu}$ of Rus01 (13.x), made explicit for the only place where it is needed, the real variables (no domain loss).}
\end{lemma}
\begin{proof}
\proofstep{Global flow}\label{stepP:p11a:rl-global} $\varphi_1$ is smooth, bounded and Lipschitz, so solutions are global; periodicity of $\varphi_1$ gives $\Theta^t(\theta+k)=\Theta^t(\theta)+k$ (uniqueness); $\lvert\Theta^t-\theta\rvert\le\lvert t\rvert\sup\lvert\varphi_1\rvert$.

\proofstep{Derivative}\label{stepP:p11a:rl-deriv} $Y(t)=D\Theta^t(\theta)$ solves $Y'=D\varphi_1(\Theta^t\theta)Y$, $Y(0)=\mathbf 1$ (\ref{bpP:p11a:x-ode}); $\lvert Y(t)\rvert\le e^{\Lambda\lvert t\rvert}$ and $\lvert Y(t)-\mathbf 1\rvert\le\int_0^{\lvert t\rvert}\Lambda e^{\Lambda\tau}\,d\tau=e^{\Lambda\lvert t\rvert}-1$.

\proofstep{The constant}\label{stepP:p11a:rl-const} The $i$-th row sum of $D\varphi_1(x)$ is $\sum_j\lvert\partial_j\varphi_{1,i}(x)\rvert\le\sum_{k}\lvert\hat\varphi_{1,i,k}\rvert\,2\pi\sum_j\lvert k_j\rvert\le2\pi T\sum_{0<\lvert k\rvert_1\le T}\lvert\hat\varphi_{1,i,k}\rvert$ ($\lvert\chi_k\rvert=1$ on $\R^N$, $k=0$ drops out). With the decay bound, the sum is $\le M_1\Lam{r_*}/(\pi D)$.
\end{proof}

%%%%%%%%%%%%%%%%%%%%%%%%%%%%%%%%%%%%%%%%%%%%%%%%%%%%%%%%%%%%%%%%%%%%%%%%%%%%%%%%
\subsection{The canonical step}\label{bpP:p11a:sec-step}

\begin{proposition}[Canonical step in the maximal case]\label{bpP:p11a:canonical-step}
\textbf{Data.} $N\ge1$; $E\subset\R^N$, $t>0$, $B=E+t$; $r,s>0$; $h:B\to\C$ and $\omega:B\to\C^N$ holomorphic, real on real points; $R:\Dd rsB\to\C$ admissible with $\varepsilon=\nrm R{r,s;B}<\infty$ (real on real points, if reality is wanted); $\varepsilon=0$ is allowed, and then $A_1=0$ below, $R_*=0$, $\varphi=0$, $\Phi=\mathrm{id}$, $R_+=0$ (the argument never divides by $\varepsilon$ or $A_1$; this is Lemma~\ref{bpP:p11a:affine-flow} with $A_1=0$; review P11A2-03); $N_b(\theta,\eta)=h(b)+\omega(b)\cdot\eta$ and $H=N_b+R$.

\textbf{Step parameters.} $T>0$; $d,\delta,\sigma>0$ with $r_+:=r-d-\delta-2\sigma>0$; $q\in(0,\frac12]$; $D>0$. Put $E_T=\varepsilon_N(T,d)$ and $A_1=\Lam\delta(1+E_T)\varepsilon/(\pi Ds)$.

\textbf{Hypotheses.}
\begin{description}
\item[(C1)] (cutoff error) $E_T\le1$ (for example $E_T\le\alpha\le1$ by Lemma~\ref{bpP:p11a:loss-choice} with $d=d_T(\alpha)\le1$);
\item[(C2)] (divisor) $\lVert D\omega\rVert\le L$ on $B$, $tTL\le D/2$, and $\lvert k\cdot\omega(b')\rvert\ge D$ for $b'\in E$, $0<\lvert k\rvert_1\le T$ (for $E=K_\nu$, $D=\gamma T^{-\tau_D}$ this is Lemma~\ref{bpP:p11a:divisor-wide});
\item[(C3)] (smallness) $4(N+1)A_1\le\sigma$, i.e.\ $4(N+1)\Lam\delta(1+E_T)\varepsilon\le\pi Ds\sigma$.
\end{description}
\textbf{Conclusions.} Let $R_*=R_{*0}+R_{*1}\cdot\eta$ be the cutoff of Proposition~\ref{bpP:p11a:cutoff} (with $T,d,q$), $\Delta h=\langle R_{*0}\rangle$, $\Delta\omega=\langle R_{*1}\rangle$, $N_+=h+\Delta h+(\omega+\Delta\omega)\cdot\eta$, $D_+=\Dd{r_+}{qs}B$, and $\varphi=\varphi_0+\varphi_1\cdot\eta$ the solution of Proposition~\ref{bpP:p11a:homological-maximal} with $P=R_*$, $r_*=r-d$, $M_0=(1+E_T)\varepsilon$, $M_1=M_0/s$. Then:
\begin{enumerate}[label=(\roman*)]
\item (normal form) $\Delta h,\Delta\omega$ are holomorphic on $B$, $\lvert\Delta h\rvert\le(1+E_T)\varepsilon\le2\varepsilon$, $\lvert\Delta\omega\rvert_\infty\le(1+E_T)\varepsilon/s\le2\varepsilon/s$; $N_+$ has the form of $N$.
\item (the map) $\Phi_b=X^1_{\varphi_b}$ is defined and holomorphic on $D_+$ jointly in $(\theta,\eta,b)$, with values in $\Dd{r_++\sigma}{s}B$; precisely $\lvert\Phi_\theta-\theta\rvert_\infty\le A_1\le\sigma/(4(N+1))$, $\lvert\Phi_\eta-\eta\rvert_\infty\le A_\eta:=(N+1)sA_1/\sigma\le s/4$; $\Phi$ is simple (Lemma~\ref{bpP:p11a:affine-flow}(c)), canonical, $\Phi(\theta+k,\eta)=\Phi(\theta,\eta)+(k,0)$, real at real points if the data are real; on $D_+$: $\lvert D\Phi-\mathbf 1\rvert_\nu\le2\kappa$ with $\kappa=2(N+1)A_1/\sigma\le\frac12$ and $\nu=\nu_{s/2,\sigma}$; $\Phi$ is invertible with inverse $X^{-1}_\varphi$ (Lemma~\ref{bpP:p11a:affine-flow}(e)).
\item (conjugation) $H_b\circ\Phi_b=N_++R_+$ on $D_+$, where $R_+$ is holomorphic on $D_+$, $1$-periodic in $\theta$, real at real points if the data are real, and
\begin{multline*}\nrm{R_+}{D_+}\ \le\ \varepsilon\Bigl(\frac{q^2}{1-q}+2E_T\Bigr)+\frac{N(4N+5)(5+4E_T)(1+E_T)\,\Lam\delta}{\pi\,D\,s\,\sigma}\,\varepsilon^2\\ \le\ \varepsilon\Bigl(2q^2+2E_T\Bigr)+\frac{18\,N(4N+5)\Lam\delta}{\pi Ds\sigma}\varepsilon^2 .\end{multline*}
\item (real Lipschitz) for real $b\in E$ and real $\theta$: $\lvert D\Theta_b-\mathbf 1\rvert_\infty\le e^{\Lambda_\mathbb R}-1$, with $\Lambda_\mathbb R=2T(1+E_T)\varepsilon\,\Lam{r-d}/(Ds)$, and $\Theta_b(\theta+k)=\Theta_b(\theta)+k$.
\item (parameter derivatives) for $0<t'<t$, $j\ge1$, on $E+t'$ and at every $(\theta,\eta)\in D_+$: $\lVert D_b^j\Delta h\rVert\le j^j\frac{(1+E_T)\varepsilon}{(t-t')^j}$, $\lVert D_b^j\Delta\omega\rVert\le j^j\frac{(1+E_T)\varepsilon}{s(t-t')^j}$, $\lVert D_b^j(\Phi_\theta-\theta)\rVert\le\frac{j^jA_1}{(t-t')^j}$, $\lVert D_b^j(\Phi_\eta-\eta)\rVert\le\frac{j^jA_\eta}{(t-t')^j}$, $\lVert D_b^jR_+\rVert\le\frac{j^j\nrm{R_+}{D_+}}{(t-t')^j}$.
\end{enumerate}
\source{Rus01 \S\S13.6--13.12 pp.~159--164 (the estimate of the remainder after one step; not read: only the structure is used, the constants are recomputed here), Lemma~11.1 p.~148 (Poisson bracket estimate) replaced by Cauchy estimates, \S12; Pos93 pp.~211--212 through \Nref{bpN:p5:taylor}. Constants are ours (D-9).}
\end{proposition}
\begin{proof}
\proofstep{Cutoff}\label{stepP:p11a:cs-cutoff} Proposition~\ref{bpP:p11a:cutoff} with the given $T,d,q$ (here $d<r$ since $r_+>0$): $R_*\in$ trigonometric polynomials of degree $\le T$, with $\nrm{R_{*0}}{r-d;B}\le M_0$, $\nrm{R_{*1,j}}{r-d;B}\le M_1$, $\nrm{R_*}{r-d,s;B}\le2M_0$, $\nrm{R-R_*}{r-d,qs;B}\le\varepsilon(\frac{q^2}{1-q}+2E_T)$.

\proofstep{Homological equation}\label{stepP:p11a:cs-hom} Proposition~\ref{bpP:p11a:homological-maximal} with $r_*=r-d$ (and $\delta<r-d$), $P_0=R_{*0}$, $P_1=R_{*1}$: hypothesis (C2). We get $\varphi$ on $T_{r_1}$, $r_1=r-d-\delta$, with $\nrm{\varphi_{1,j}}{r_1;B}\le\Lam\delta M_1/(\pi D)=A_1$ and $\nrm{\varphi_0}{r_1;B}\le\Lam\delta M_0/(\pi D)=sA_1$, and (i) from part (b) of that proposition ($\Delta h=\langle R_{*0}\rangle$, $\lvert\Delta h\rvert\le M_0$, etc.). Moreover $\{N_b,\varphi\}=-(R_*-\langle R_*\rangle)=:w$.

\proofstep{Flow}\label{stepP:p11a:cs-flow} Lemma~\ref{bpP:p11a:affine-flow} with $r_1=r-d-\delta=r_++2\sigma$ (so $r_1-2\sigma=r_+$), $A_1$ as above, hypothesis (C3) $=4(N+1)A_1\le\sigma$, fibrewise in $b\in B$ and with the parameter statement (f). The domain $\Omega_{\mathrm f}=T_{r_1-2\sigma}\times B_{s/2}=T_{r_+}\times B_{s/2}\supset D_+$ fibrewise ($q\le\frac12$), so (ii) follows from Lemma~\ref{bpP:p11a:affine-flow}(b)--(f); the displacements are the stated $A_1,A_\eta$.

\proofstep{Taylor identity}\label{stepP:p11a:cs-taylor} Fix $b\in B$ and $z\in D_+\cap\{b\text{ fixed}\}$. Put $\Omega=T_{r_1-\sigma}\times B_s$; $N_b$ is entire and $R$ is holomorphic on $T_r\times B_s\supset\Omega$ ($r_1\le r$); the trajectory $X^t_\varphi(z)$ stays in $\Omega$ (Lemma~\ref{bpP:p11a:affine-flow}(b)) and $z\in\Omega^{1/2}$. \Nref{bpN:p5:taylor} (the displayed formula for $(h+f_0)\circ X^1$) with $h=N_b$, $f_0=R$, gives
\[H\circ\Phi(z)=N_b(z)+\{N_b,\varphi\}(z)+R(z)+\mathcal R(z),\qquad \mathcal R(z)=\int_0^1\{u_t,\varphi\}(X^t_\varphi z)\,dt,\quad u_t=(1-t)w+R.\]
With $\{N_b,\varphi\}=w=-(R_*-\langle R_*\rangle)$: $N_b+w+R=N_b+\langle R_*\rangle+(R-R_*)=N_++(R-R_*)$. Hence $H\circ\Phi=N_++R_+$ with $R_+=(R-R_*)+\mathcal R$ on $D_+$. Holomorphy of $R_+$ in $(\theta,\eta,b)$ is that of $H$ and of $\Phi$ (joint, Lemma~\ref{bpP:p11a:affine-flow}(f)); periodicity and reality follow from those of $H,\Phi$ and $N_+$.

\proofstep{Size of $u_t$}\label{stepP:p11a:cs-u} $u_t$ is holomorphic on $T_{r-d}\times B_s$. There $\lvert w\rvert\le\lvert R_*\rvert+\lvert\langle R_*\rangle\rvert\le4(1+E_T)\varepsilon$ ($\langle R_*\rangle(\eta)$ is an average of values of $R_*$ at real $\theta$, hence bounded by $\nrm{R_*}{r-d,s;B}\le2(1+E_T)\varepsilon$), so $\lvert u_t\rvert\le U:=\varepsilon(5+4E_T)\ (\le9\varepsilon)$ for $t\in[0,1]$.

\proofstep{Cauchy estimates for the bracket}\label{stepP:p11a:cs-bracket} Let $z\in D_+$ and $z_t=X^t_\varphi(z)=(\theta_t,\eta_t)$. By the displacement bounds, $\theta_t\in T_{r_++A_1}$, $\lvert\eta_t\rvert<qs+A_\eta\le qs+s/4$. The polydisc around $z_t$ of radii $(\sigma,\rho_u)$ with $\rho_u=(1-q)s/2\ \ (\ge s/4)$ lies in $T_{r-d}\times B_s$: $r_++A_1+\sigma\le r_++\frac98\sigma\le r-d$, and $qs+\frac s4+\frac{(1-q)s}2\le s$ for $q\le\frac12$. By \Kref{bp:p1:ext-cauchy} (one variable at a time), $\lvert\partial_{\theta_i}u_t(z_t)\rvert\le U/\sigma$ and $\lvert\partial_{\eta_i}u_t(z_t)\rvert\le U/\rho_u\le4U/s$. Also $z_t\in\Omega$ so $\lvert\partial_{\eta_i}\varphi\rvert\le A_1$, $\lvert\partial_{\theta_i}\varphi\rvert\le A_\eta$ at $z_t$ (Lemma~\ref{bpP:p11a:affine-flow}(a)). Hence
\[\lvert\{u_t,\varphi\}(z_t)\rvert\le\sum_{i=1}^N\bigl(\lvert\partial_{\theta_i}u_t\rvert\lvert\partial_{\eta_i}\varphi\rvert+\lvert\partial_{\eta_i}u_t\rvert\lvert\partial_{\theta_i}\varphi\rvert\bigr)\le N\Bigl(\frac{UA_1}\sigma+\frac{4U}s\cdot\frac{(N+1)sA_1}\sigma\Bigr)=\frac{N(4N+5)UA_1}\sigma.\]

\proofstep{Assembly}\label{stepP:p11a:cs-assemble} Integrate over $t\in[0,1]$: $\lvert\mathcal R\rvert\le N(4N+5)UA_1/\sigma$ on $D_+$; insert $U=\varepsilon(5+4E_T)$, $A_1=\Lam\delta(1+E_T)\varepsilon/(\pi Ds)$ and the bound of $R-R_*$ (Step~\ref{stepP:p11a:cs-cutoff}; $D_+\subset T_{r-d}\times B_{qs}$). The second form uses $q^2/(1-q)\le2q^2$, $E_T\le1$, $(5+4E_T)(1+E_T)\le18$.

\proofstep{Real Lipschitz bound}\label{stepP:p11a:cs-lip} Lemma~\ref{bpP:p11a:real-lipschitz} with $r_*=r-d$, $M_1=(1+E_T)\varepsilon/s$ and Proposition~\ref{bpP:p11a:homological-maximal}(e); for real $b\in E$, $\varphi_1$ is real (Prop.~\ref{bpP:p11a:homological-maximal}(d)) and the real flow of $\dot\theta=\varphi_1$ is the restriction of $\Theta_b$ (uniqueness).

\proofstep{Parameter derivatives}\label{stepP:p11a:cs-param} Lemma~\ref{bpP:p11a:param-cauchy} applied, for fixed $(\theta,\eta)\in D_+$, to $b\mapsto\Delta h,\Delta\omega,\Phi_\theta-\theta,\Phi_\eta-\eta,R_+$ on $B=E+t$, using their bounds from (i), (ii), (iii).
\end{proof}

\begin{remark}[Reading the statement; the next step]\label{bpP:p11a:step-reading}
(a) \emph{Domain bookkeeping.} The angular strip loses $d+\delta+2\sigma$ ($d=d_T$ Fourier, $\delta$ homological, $2\sigma$ Cauchy/flow); the action radius becomes $qs$ ($q\le\frac12$); the parameter domain $B=E+t$ is \emph{not} reduced by the step itself (the parameter Cauchy estimates (v) cost $t-t'$ only if derivatives are wanted). The output $(N_+,R_+)$ on $D_+$ has the input form with $(r,s,\varepsilon,h,\omega)\to(r_+,qs,\varepsilon_+,h+\Delta h,\omega+\Delta\omega)$, $\varepsilon_+\le$ the bound of (iii). \emph{(b) What the next step needs on the frequency:} holomorphy of $\omega+\Delta\omega$ on $B$ and the Lipschitz bound $L_+\le L+\sup\lVert D\Delta\omega\rVert$, with $\lVert D\Delta\omega\rVert\le(1+E_T)\varepsilon/(s(t-t'))$ on $E+t'$ (v) ($j=1$); so the next step runs on $E_+\subset E$ with $t_+\le t'$: File~11a-3 sets $t_{\nu+1}=t'$ and checks $t_{\nu+1}T_{\nu+1}L_{\nu+1}\le D_{\nu+1}/2$. \emph{(c)} The only place where the \emph{small divisor} enters is $D$ in (C2)--(C3), (iii), (iv); the cutoff $T$ enters through $E_T$, $d$ and (C2). \emph{(d)} The first term of (iii) is linear in $\varepsilon$ (cutoff errors $q^2$, $E_T$); only the second is quadratic. In Rus01's scheme $\varepsilon_{\nu+1}=\varepsilon_\nu\delta^{\kappa}$ is only geometric, so one needs $2q_\nu^2+2\alpha_\nu\lesssim\delta^{\kappa}/2$ and the quadratic term $\lesssim\varepsilon_\nu\cdot\varepsilon_\nu/(D_\nu s_\nu\sigma_\nu)\ll\delta^\kappa\varepsilon_\nu$ (File~11a-3).
\end{remark}

%%%%%%%%%%%%%%%%%%%%%%%%%%%%%%%%%%%%%%%%%%%%%%%%%%%%%%%%%%%%%%%%%%%%%%%%%%%%%%%%
\subsection{Sources, deviations, errata, what File~11a-3 takes}\label{bpP:p11a:notes2}

\paragraph{Sources read.} Rus01 \S\S7.1--7.5 (pp.~134--136), \S8.1--8.3 (p.~137), \S\S11.1--11.4 (pp.~148--153), \S12 (pp.~154--155), \S13.1--13.2 (p.~156; schedule only); Rus05 \S\S9--11 (pp.~26--31) as comparison (equation (46)); the review \texttt{reviews/p11a-russmann-check.md} (rows \texttt{analytic-cutoff}, \texttt{homological-maximal}, \texttt{canonical-step}; errata R05-1, R05-2) and \texttt{reviews/p11a1-review-1.md} (P11A1-P02); track~N Parts~4 and~5. Not read: Rus01 \S10 (pp.~143--147, the flow of the full vector field with elliptic/hyperbolic variables), \S\S13.3--13.12 (pp.~157--164): the maximal case needs neither, the one-step remainder is recomputed above.

\paragraph{Deviations (all recorded; D-numbers as in \texttt{PLAN-11a.md}).}
\begin{itemize}
\item \textbf{D-4} (Fourier cutoff): Lemma~\ref{bpP:p11a:fourier-cutoff} replaces Rus01 Thm~7.3 and Rus05 by a short proof from Cauchy decay; no multipliers; truncation in $\lvert k\rvert_1$; the constant $c_N$ is explicit; the loss is $d_T$ of Lemma~\ref{bpP:p11a:loss-choice}, not Rus01's $\rho=\frac nT\log\frac{4n}\alpha$ with error $\frac45\alpha$. The cube/Euclid rescaling that causes R05-2 does not occur: $T$ and $d$ are used in the same $\ell^1$ lattice norm throughout.
\item \textbf{D-5} (algebraic cutoff): Lemma~\ref{bpP:p11a:taylor-cutoff}, Taylor remainder of order $2$ with ratio $q$ (Rus01's $\delta^{\beta}$); no appeal to [18] Lemma~6; only degree $\le1$ is needed.
\item \textbf{D-10} (sup norms): all estimates are in sup norms on $\Dd rsB$ (as Rus01 \S7--\S13), not in the weighted norm of track~N; the passage between them is not needed, since only $\cut_T$ and the homological equation use Fourier coefficients (Cauchy decay, \Kref{bp:p1:decay}).
\item \textbf{D-11} (homological equation): solved coefficientwise, Rus01 (11.12) with $p=q=0$; the loss $\delta$ enters through $\Lam\delta=\coth(\pi\delta)^N-1$, not through $(2T+1)^n$ of (11.21); the divisor is $D/2$ on all of $B$ (extension from $E$, Lemma~\ref{bpP:p11a:divisor-perturb}).
\item \textbf{D-12} (sign): $X_\phi=(\partial_\eta\phi,-\partial_\theta\phi)$ and $\omega\cdot\partial_\theta\varphi=+(R_*-\langle R_*\rangle)$ (track~N's convention). This is the same convention as Rus01's $\{S,N\}=F-[F]$, (11.12) after row/column conversion (Rus01 \S10.1.8 uses $J_{\mathrm{src}}=\bigl(\begin{smallmatrix}0&-\mathbf 1\\ \mathbf 1&0\end{smallmatrix}\bigr)$ acting on row vectors, so $S_xJ_{\mathrm{src}}=(S_y,-S_x)$; the displayed matrices differ only because of row versus column multiplication); the positive homological sign corrects the negative sign $\partial_\omega\varphi=-(R-\langle R\rangle)$ of \texttt{PLAN-11a.md}.
\item \textbf{D-13} (higher parameter Cauchy): by polydisc iteration, Lemma~\ref{bpP:p11a:param-cauchy}; no polarization constant.
\item \textbf{D-14} (flow): Rus01 Thm~8.1 (flow in $\C^m$ with $\lvert f(x)\rvert\le Me^{\lvert x\rvert}$ and the Poisson bracket estimate 11.1) is replaced by \Nref{bpN:p5:flow} and Cauchy estimates of $\{u,\varphi\}$, which use the explicit polydisc geometry of the strips.
\item \textbf{D-15} (strip bookkeeping): the angular loss is $d+\delta+2\sigma$ per step; it is a datum of File~11a-3 (condition (I0), $\sum_\nu(d_\nu+\delta_\nu+2\sigma_\nu)\le s/2$), with $\sum_\nu d_\nu$ given by Lemma~\ref{bpP:p11a:loss-sum}.
\end{itemize}

\paragraph{Suspected errata.} (1) Rus05 p.~30, equation (46) and p.~31 (R05-2, from the review): the dimension rescaling $T\to T/\sqrt n$ produces $E(T\rho/n)$, not $E(T\rho)$; not used here, recorded in the ledger by the orchestrator. (2) No new defect found in Rus01 pp.~134--136, 148--155 beyond typographical notation (the convention of rows and columns in \S12). The statement of Rus01 Thm~7.3 is correct as quoted, but its proof imports [21] (a theorem of Akhiezer proved elsewhere); see D-4.

\paragraph{Doubts about frozen statements.} None: this file touches no frozen declaration.

\paragraph{What File~11a-3 must take from this file.}
\begin{enumerate}[label=(T\arabic*)]
\item \emph{Per step, Proposition~\ref{bpP:p11a:canonical-step}} with a stage index: parameters $(T_\nu,d_\nu,\delta_\nu,\sigma_\nu,q_\nu,D_\nu)$; the hypotheses (C1)--(C3) become the schedule inequalities: $d_\nu=d_{T_\nu}(\alpha_\nu)\le1$ with $\alpha_\nu\le1$ (Lemma~\ref{bpP:p11a:loss-choice}); $D_\nu=\gamma T_\nu^{-\tau_D}$ with the divisor hypothesis from Lemma~\ref{bpP:p11a:divisor-wide} (here $E=K_{\nu+1}$, the retained set, with the divisor bound of Lemma~\ref{bpP:p11a:finite-divisor}; $T=T_\nu$, $\tau=\tau_D\ge0$); $t_\nu T_\nu L_\nu\le D_\nu/2$; and $4(N+1)\Lam{\delta_\nu}(1+E_{T_\nu})\varepsilon_\nu\le\pi D_\nu s_\nu\sigma_\nu$.
\item \emph{The recursions:} $r_{\nu+1}=r_\nu-d_\nu-\delta_\nu-2\sigma_\nu$; $s_{\nu+1}=q_\nu s_\nu$; $\varepsilon_{\nu+1}\le\varepsilon_\nu(2q_\nu^2+2\alpha_\nu)+\frac{18N(4N+5)\Lam{\delta_\nu}}{\pi D_\nu s_\nu\sigma_\nu}\varepsilon_\nu^2$; $\omega_{\nu+1}=\omega_\nu+\Delta\omega_\nu$ with $\lvert\Delta\omega_\nu\rvert_\infty\le2\varepsilon_\nu/s_\nu$ (Rus01 $3\varepsilon_\nu/(2s_\nu)$ when $E_T\le\frac12$); $h_{\nu+1}=h_\nu+\Delta h_\nu$, $\lvert\Delta h_\nu\rvert\le2\varepsilon_\nu$.
\item \emph{Displacement and Lipschitz data for File~11a-4:} $\lvert X_{\nu,\theta}-\theta\rvert_\infty\le A_{1,\nu}$, $\lvert X_{\nu,\eta}-\eta\rvert_\infty\le A_{\eta,\nu}$ on $D_{+,\nu}$; $\lvert DX_\nu-\mathbf 1\rvert_\nu\le2\kappa_\nu$ on the complex domain in the stage-dependent weighted norm $\nu=\nu_{s_\nu/2,\sigma_\nu}$ (the summable quantity $\sum\kappa_\nu$ is $2(N+1)\sum A_{1,\nu}/\sigma_\nu$), \emph{and} the unweighted bound $\lvert DX_\nu-\mathbf 1\rvert_\infty\le2\kappa_\nu+\kappa_\nu s_\nu/\sigma_\nu$ in the coordinates $(\theta,\eta)$ (diagonal blocks $\le2\kappa_\nu$, lower-left block $\le\kappa_\nu s_\nu/\sigma_\nu$ by Lemma~\ref{bpP:p11a:affine-flow}(d); maximal row sum), whose sum over $\nu$ the consumer must control (File~11a-3, Lemma~\ref{bpP:p11a:exported-sums}(E2): $\le2^{-12}$); the weighted bound alone is not a bound for the composition in a fixed norm (review P11A2-04); and on the reals $\lvert D\Theta_\nu-\mathbf 1\rvert_\infty\le e^{\Lambda_{\mathbb R,\nu}}-1$ with $\Lambda_{\mathbb R,\nu}=2T_\nu(1+E_{T_\nu})\varepsilon_\nu\Lam{r_\nu-d_\nu}/(D_\nu s_\nu)$ (the analogue of Rus01's $e^{cC_\nu T_\nu}$). The factor $\Lam{r_\nu-d_\nu}$ is large when $r_\nu-d_\nu$ is small; File~11a-3 keeps $r_\nu\ge s/2$-type lower bounds on the strip.
\item \emph{Parameter regularity (for the chain extension, deviation D-7):} (v) of the proposition gives $\lVert D_b^j\Delta\omega_\nu\rVert\le j^j(1+E_T)\varepsilon_\nu/(s_\nu(t_\nu-t_{\nu+1})^j)$ on $E+t_{\nu+1}$, for $j\le\mu+1$: this is the \emph{holomorphic} size of the increments, from which the mollified cut-offs of File~11a-3 produce the $C^\mu$-size.
\item \emph{The interface to File~11a-1:} Lemma~\ref{bpP:p11a:divisor-wide} is the only point where $\tilde\omega_k$, $\nu(k)$, (I1)--(I3) enter this file; it turns them into the single divisor $D_\nu=\gamma T_\nu^{-\tau_D}$ used by the homological equation, on $E$ only; the finite-stage invariant that pays its hypothesis on the retained set is Lemma~\ref{bpP:p11a:finite-divisor} (File~11a-3), which supersedes the limit-based form of \ref{bpP:p11a:I3} (Remark~\ref{bpP:p11a:i3-corrected}; review P11A2-01); the passage to $E+t$ is Lemma~\ref{bpP:p11a:divisor-perturb} (requires $t_\nu T_\nu L_\nu\le D_\nu/2$, another schedule inequality).
\end{enumerate}

\paragraph{What is reused and what is new.} Track~N: \Nref{bpN:p5:flow}, \Nref{bpN:p5:symplectic}, \Nref{bpN:p5:taylor}, \Nref{bpN:p5:setting}, \Nref{bpN:p4:ext-param} (as cited). Not used: \Nref{bpN:p4:projections} (weighted-norm form of the cutoff, $e^{-2\pi K\sigma}$, which has no polynomial prefactor but requires a weighted norm of $R$, i.e.\ a $\coth^N$ conversion at every step, \Nref{bpN:p4:sup-weighted}; the sup-norm version above is cheaper here because the pairing of $\lvert k\rvert_1\le T$ with a width $d\sim T^{-1}\log T$ is the only loss), \Nref{bpN:p5:composition}, \Nref{bpN:p5:transform}, \Nref{bpN:p5:canonical-map} (weighted norms; Part~6's Iterative Lemma is not usable, see \texttt{PLAN-11a.md}). Track~K: \Kref{bp:p1:decay}, \Kref{bp:p1:summation}, \Kref{bp:p1:exp-sums}, \Kref{bp:p1:cauchy-strip}, \Kref{bp:p1:fourier-basic}, \Kref{bp:p1:ext-cauchy}, \Kref{bp:p1:ext-mvt}, \Kref{bp:p4:lattice-count}. New in this file: all of the statements above; in particular the simple-form and real-Lipschitz statements, the loss choice and its sum, and the explicit constants of the one-step remainder.

\paragraph{External facts used.} \ref{bpP:p11a:x-cauchy}--\ref{bpP:p11a:x-ode}; \Kref{bp:p1:ext-cauchy}, \Kref{bp:p1:ext-mvt}, \Kref{bp:p1:ext-integral}; \Nref{bpN:p4:ext-param}. The multivariable Cauchy coefficient estimate is stated in the form used; the Mathlib names are likely, not verified.

\paragraph{Numerical checks.} \texttt{.drafts/P/bp-scratch-bpP-p11a-2/chk1.py}: for $N\le6$, $d\in[0.01,5]$, $T\le100$, the exact lattice tail $\sum_{\lvert k\rvert_1>T}e^{-2\pi\lvert k\rvert_1d}$ is at most $0.12\,c_N(1+d^{-N})e^{-\pi Td}$ and $\varepsilon_N(T,d_T)\le0.34\,\alpha$ whenever $d_T\le1$ (evidence only). \texttt{chk2.py}: $N=1$, $R=\varepsilon_0\cos(2\pi\theta+\eta)$, $\omega=\sqrt2$: the real remainder $H\circ\Phi-N_+$ computed by numerical flow is $3\cdot10^{-5}$ ($\approx\varepsilon_0q^2/2$, the Taylor part; a wrong sign of $\varphi$ would give $\approx\varepsilon_0$), below the bound of (iii) ($3.8\cdot10^{-2}$).

\paragraph{Forward references and label index.} No forward reference to an unwritten label. Labels: \texttt{c2-conv}, \texttt{domain-def}, \texttt{fourier-cutoff}, \texttt{loss-choice}, \texttt{loss-sum}, \texttt{taylor-cutoff}, \texttt{cutoff}, \texttt{param-cauchy}, \texttt{divisor-wide}, \texttt{divisor-perturb}, \texttt{homological-maximal}, \texttt{affine-flow}, \texttt{real-lipschitz}, \texttt{canonical-step}, \texttt{step-reading}.
\par\endgroup
